%=============================================================================!
% 18.2  BASIS FUNCTIONS AND OVERFITTING                                       !
%=============================================================================!
\section{Basis functions and overfitting}
\label{sec:lo-basis}

A curve drawn freehand through a dozen noisy points can be made to pass
through every one of them, if it is allowed to wiggle; between the
points it then swings far from any sensible value, and a new point
measured in a gap lands nowhere near it. The curve has learned the
noise. A smoother curve that passes near the points, rather than
through them, predicts the gaps better. Choosing how freely a model may
bend is the first decision of any fit.

\subsection*{A curve through noisy readings}

The readings here are argon's pair energy $\varphi(r)$ between
\SI{3.3}{\angstrom} and $2\sigma = \SI{6.8}{\angstrom}$, where the switched
potential of Chapter~12 is the plain Lennard-Jones energy, read at 15
distances drawn at random with an error of \SI{0.2}{\milli\electronvolt}
added, as a first-principles calculation adds its own small errors
(\cref{fig:lo-basis}a). The model is a polynomial,
\begin{equation*}
   \hat\varphi(r) = \sum_{j=0}^{p-1}w_js^j , \qquad
   s = \frac{r - \SI{5.05}{\angstrom}}{\SI{1.75}{\angstrom}} ,
\end{equation*}
in the scaled distance $s$, which runs from $-1$ to 1 over the range; the
functions $s^j$ are its \emph{basis functions}, and the columns of the
design matrix are their values at the readings. The scaling matters for
the arithmetic, not for the fit: the powers of $r$ itself range from 1
to $6.8^{14} \approx \num{5e11}$, and \cref{sec:lo-descent} measures what
that does. The weights come from the normal equations
(\cref{eq:lo-normal}), and the fit is judged on 200 distances spread
evenly over the range, without noise, which the fit never saw: the
\emph{test}.

\Cref{fig:lo-basis}(b) tells the story. As the degree of the polynomial
rises from 0, the rms miss on the 15 training readings falls, from
\num{3.38} to \SI{0.20}{\milli\electronvolt} at degree 8, near the
noise, and then further, to zero at degree 14, where 15 weights pass a
curve through 15 readings exactly. The miss on the test falls with it at
first, to \SI{0.28}{\milli\electronvolt} at degree 8, and then rises
steeply: \SI{3.9}{\milli\electronvolt} at degree 9, 140 at 12 and
\SI{41}{\electronvolt} at 14, where the curve swings between the readings
(\cref{fig:lo-basis}a). A model with too few weights misses the curve,
\emph{underfitting}; one with too many learns the noise of its readings
as if it were part of the curve, \emph{overfitting}, and the training
miss falling below the noise of the readings is its sign. Only data the
model has not seen can tell the two apart, since the training miss
improves with every weight added.

The functions themselves can carry knowledge. The Lennard-Jones energy is
a sum of $(\sigma/r)^{12}$ and $(\sigma/r)^6$, and a model built on those
two functions has two weights to find. Fitted to the same 15 readings,
they come out as \num{0.04051} and \SI{-0.04073}{\electronvolt}, against
$4\varepsilon = \SI{0.04136}{\electronvolt}$ and $-4\varepsilon$, and
the test miss is \SI{0.090}{\milli\electronvolt}, a third of the best
polynomial's with two weights against nine. A learned potential faces
the same choice on a larger scale: a model that knows the forces between
atoms depend only on their distances and on the elements needs far fewer
readings than one that must discover it, which is what planned Chapter~20 builds
into the description of an atom's surroundings.

\begin{figure}[htb]
\centering
\includegraphics{ch18_learning/basis}
\caption{Basis functions and overfitting. (a) Argon's pair energy (grey
dashed) read at 15 distances with an error of
\SI{0.2}{\milli\electronvolt} (dots), and the polynomials of degree 3, 8
and 14 in the scaled distance fitted to the readings by least squares.
(b) The rms miss on the training readings and on a test of 200
noise-free distances against the degree, with the noise (dotted) and the
test miss of the two functions $(\sigma/r)^{12}$ and $(\sigma/r)^6$
(cross). At degree 14 the curve passes through every reading, and its
training miss, zero, is drawn on the axis.}
\label{fig:lo-basis}
\end{figure}

\begin{keyidea}
A model linear in its weights is a sum of basis functions, and the
number of them sets how freely it bends. Too few underfit; too many
learn the noise, which shows as a training miss below the noise of the
readings and a test miss that grows. Only readings the model has not
seen judge it. Functions that carry the physics need fewer weights and
fewer readings.
\end{keyidea}

\notebook{18}{change the degree of the polynomial and watch the training
and test misses part}

\begin{selfcheck}
A polynomial of degree 14 is fitted to 15 readings. Why is its training
miss zero whatever the readings, and what does that say about the
training miss as a measure of a model?
\end{selfcheck}
